\section*{الفصل \ref{ch:g9:alternating-current} --- التيار المتناوب}

\begin{solution}{exo:g9:alternating-current:1}
المستمرّ يسري دائمًا في جهة واحدة تحت \omterm{def:g8:voltage-voltmeter:voltage}{توتر} ثابت --- نظام
\omterm{def:g3:first-electric-circuit:battery}{البطارية}. والمتناوب يعكس اتجاهه دوريًّا تحت \omterm{def:g8:voltage-voltmeter:voltage}{توتر} متأرجح ---
نظام المأخذ.
\end{solution}

\begin{solution}{exo:g9:alternating-current:2}
$T$: \omterm{def:g2:measuring-time:duration}{مدة} دورة كاملة واحدة؛ و$f$: الدورات في الثانية؛
و$f = 1/T$ و$T = 1/f$. فعند \qty{50}{Hz}: $T = \qty{0.02}{s}$.
وعند $T = \qty{0.01}{s}$: $f = \qty{100}{Hz}$.
\end{solution}

\begin{solution}{exo:g9:alternating-current:3}
مئة مرة: انعكاسان في كل دورة من الدورات الخمسين.
\end{solution}

\begin{solution}{exo:g9:alternating-current:4}
$T = 4 \times 5 = \qty{20}{ms} = \qty{0.02}{s}$؛ و$f = 1/0.02
= \qty{50}{Hz}$.
\end{solution}

\begin{solution}{exo:g9:alternating-current:5}
\qty{230}{V} هي \omterm{prop:g9:alternating-current:effective}{التوتر الفعّال} --- أي القيمة المستمرّة الثابتة
التي كانت ستسخّن بالقدر نفسه، وهي ما تقرأه مقاييس \omterm{def:g8:voltage-voltmeter:voltage}{التوتر}
المتناوب؛ و\qty{325}{V} هي القمّة، أي الذروة التي تلمسها
التأرجحة فعلًا، في كل جهة، خمسين مرة في الثانية.
\end{solution}

\begin{solution}{exo:g9:alternating-current:6}
التسخين لا يبالي بالاتجاه --- فالوشيعة تدفأ على أيّ مسير،
فتشرب الغلّايات التأرجحة مباشرة. أما الشحن فكيمياء معكوسة،
والكيمياء تقتضي اتجاهًا واحدًا ثابتًا: فالطوبة تحوّل التأرجحة
إلى ثبات أولًا.
\end{solution}

\begin{solution}{exo:g9:alternating-current:7}
\omterm{def:g3:first-electric-circuit:battery}{البطارية}: خطّ مستوٍ عند \qty{9}{V}. والمغذّي المتناوب: موجة
ناعمة تتأرجح بين $+9$ و\qty{-9}{V}، دورة لكل \qty{20}{ms}.
واختبار النظرة: مستوٍ في مقابل متموّج.
\end{solution}

\begin{solution}{exo:g9:alternating-current:8}
الآلات الدوّارة تصنع المتناوب بطبيعتها؛ ولا يغذّي المحوّلَ
إلا المتناوب، وبدونه ما كانت الكهرباء تستطيع أن تسافر البلد
\omterm{def:g8:voltage-voltmeter:voltage}{بتوتر} عالٍ وتدخل البيوت \omterm{def:g8:voltage-voltmeter:voltage}{بتوتر} مروَّض.
\end{solution}

\begin{solution}{exo:g9:alternating-current:9}
عند الدوس البطيء تنبض دورات الدينامو القليلة في الثانية
المصباحَ نبضًا مرئيًّا --- كل قمّة ومضة، وكل صفر خفوت؛ والدوس
الأسرع يضاعف النبضات حتى تدمجها العين. وصور السينما الأربع
والعشرون تحدّد معدّل الدمج التقريبي للعين: فدون نحو
\qty{25}{Hz} من نبضات \omterm{def:g1:light-and-shadows:source}{الضوء} يزعج الخفقان؛ وفوقها يُقرأ
التوهّج ثابتًا.
\end{solution}

\begin{solution}{exo:g9:alternating-current:10}
قمم عند نحو \qty{325}{V} --- ومئة منها في الثانية، خمسون
موجبة وخمسون سالبة.
\end{solution}

\begin{solution}{exo:g9:alternating-current:11}
$T = 1/60 \approx \qty{0.017}{s}$؛ والقمّة $\approx 120 \times
1.4 \approx \qty{170}{V}$. والمسخّنات والمصابيح تتكيّف بلا
مبالاة؛ والشواحن الحديثة تقرأ على لافتاتها المدى من
\qty{100}{V} إلى \qty{240}{V} والتردّدين \qty{50}{Hz}
و\qty{60}{Hz}، وتحوّل راضية --- وأمّا المعترضون
فالأجهزة القديمة ذات \omterm{def:g8:voltage-voltmeter:voltage}{التوتر} الواحد، وأيّ ساعة محرَّكة تعدّ
دورات الشبكة: فستسرع عند الستّين.
\end{solution}

\begin{solution}{exo:g9:alternating-current:12}
السطوع عند القمّتين معًا يعطي $2 \times 50 = 100$ نبضة في
الثانية. وآلة تصوير بخمس وعشرين لقطة تعاين كل \qty{0.04}{s}
--- أي أربع نبضات لكل لقطة، لكن غير مصطفّة اصطفافًا واحدًا عند
كل موضع من التعريض المتدحرج: فالفارق الطفيف بين إيقاع الآلة
وإيقاع \omterm{def:g3:first-electric-circuit:bulb}{المصباح} يرسم أشرطة تزحف عبر اللقطات المتعاقبة ---
تداخل بين ساعتين. وعند \qty{60}{Hz} ($120$ نبضة) يتغيّر
الفارق في مقابل اللقطات الخمس والعشرين، فيتغيّر عرض الأشرطة
وسرعتها --- فالنمط يفضح الشبكة المحلّية.
\end{solution}

\begin{solution}{pb:g9:alternating-current:1}
\textbf{1.} مستمرّ عند $2.25 \times 2 = \qty{4.5}{V}$:
\omterm{def:g3:first-electric-circuit:battery}{البطارية} المسطّحة الجيبية.
\textbf{2.} متناوب؛ و$T = \qty{20}{ms}$، و$f = \qty{50}{Hz}$؛
والقمّة $3 \times 2 = \qty{6}{V}$.
\textbf{3.} مستمرّ عند \qty{4.5}{V} موصول \emph{بالمقلوب} ---
فقد بُدّل السلكان: لا ضرر على الراسم، ومفيد في الدفتر.
\textbf{4.} متناوب؛ و$T = \qty{10}{ms}$، و$f = \qty{100}{Hz}$؛
والقمّة \qty{3}{V}. ونبضاته المضاعفة البالغة مئة في الثانية
تقع وراء معدّل دمج العين براحة أكبر من خفقان ب ذي الخمسين
دورة.
\textbf{5.} المقياس يبلّغ عن القيمة الفعّالة --- المكافئة
للمستمرّ في التسخين --- وهي للجيبية عند القمّة مقسومة على
$1.4$ تقريبًا: $6 \div 1.4 \approx \qty{4.3}{V}$. فالعددان
يصفان تأرجحة واحدة.
\textbf{6.} المغذّي أ وحده (بالوصل الصحيح --- وج بعد إعادة
الوصل). أما المغذّيان المتناوبان فيحتاجان إلى محوّل ---
طوبة الشاحن المقوِّمة --- بينهما وبين أيّ \omterm{def:g3:first-electric-circuit:battery}{بطارية}.
\textbf{7.} تأرجحات الشبكة الخمسون في الثانية تقع قرب إيقاع
القلب الكهربائي وقد تعطّله: فأمبيرًا \omterm{def:g8:intensity-ammeter:intensity}{بأمبير}، تعرّض التأرجحة
للخطر أكثر من المسير الثابت.
\textbf{8.} الجيبية --- التأرجحة الكاملة، ترسمها بطبيعتها
وشيعة تدور \omterm{def:g7:motion-average-speed:formula}{بسرعة} ثابتة في حقل مغناطيسي: مولّد الفصل التالي
المتناوب.
\textbf{9.} دورة لكل $4$ مربّعات (\qty{20}{ms})؛ وقمم على
$325 \div 100 = 3.25$ مربّع فوق الوسط ودونه.
\textbf{10.} $f = 1/0.02 = \qty{50}{Hz}$؛ والفعّال $\approx
325 \div 1.4 \approx \qty{230}{V}$ --- أي اللافتة، مستخرَجةً
من الشاشة.
\textbf{11.} خطّ مستوٍ عند \qty{230}{V} --- ثبات بقدرة تسخين
مساوية.
\textbf{12.} مثلًا: ``الخطّ المستوي ثبات \omterm{def:g3:first-electric-circuit:battery}{بطارية} --- اتجاه
واحد وقيمة واحدة. والموجة تأرجح الشبكة --- اتجاه وقيمة في
دورة لا تنتهي. وسعر الصرف العادل بينهما هو \omterm{prop:g9:alternating-current:effective}{التوتر الفعّال}:
\omterm{def:g5:heat-and-insulation:heat}{حرارة} واحدة، وقيمة واحدة.''
\end{solution}
